\chapter{Testing}
\label{cap:diseno_tests}

Este capítulo describe el proceso de diseño de tests basado en el conocimiento interno del sistema adquirido durante la reimplementación. El objetivo era encontrar \textbf{adversarios}: casos de test que revelen debilidades del sistema mediante análisis de caja blanca.

\section{Enfoque: Testing de Caja Blanca}

\subsection{Hipótesis}

La hipótesis central de este trabajo es que el conocimiento interno del sistema permite diseñar tests más efectivos que el testing de caja negra:

\begin{itemize}
    \item Podemos identificar \textbf{parámetros críticos} y testear cerca de sus límites
    \item Podemos entender \textbf{modos de falla} analizando la lógica del clasificador
    \item Podemos diseñar \textbf{adversarios} que exploten las características del sistema
\end{itemize}

\subsection{Conocimiento Disponible}

Del proceso de reimplementación, teníamos conocimiento detallado sobre:

\begin{itemize}
    \item \textbf{Arquitectura del recognizer:} 3 clasificadores CNN (vertical, horizontal, quad)
    \item \textbf{Clases de salida:} BLACK, RED, YELLOW, GREEN
    \item \textbf{Normalización:} means BGR específicos, scale factor 0.01
    \item \textbf{Threshold de decisión:} probabilidad $>$ 0.5 para clasificar
\end{itemize}

\subsection{Pregunta de Investigación}

La pregunta que guió el diseño del test fue:

\begin{quote}
\textit{¿Cuál es el punto exacto donde el sistema comienza a confundir el color verde con amarillo?}
\end{quote}

Esta pregunta surge del conocimiento de que verde y amarillo son colores adyacentes en el espacio de color, y que el clasificador debe tener un límite de decisión entre ellos.

\section{Test de Transición Verde → Amarillo}

\subsection{Diseño del Test}

\textbf{Idea:} Modificar gradualmente el color de un semáforo verde hacia amarillo hasta encontrar el punto donde el sistema cambia su clasificación.

\textbf{Metodología:}

\begin{enumerate}
    \item Partir de un video de semáforo en verde
    \item Aplicar modificaciones progresivas de color usando edición de video
    \item Ejecutar el sistema sobre cada versión
    \item Registrar cuándo cambia la clasificación
\end{enumerate}

\subsection{Implementación}

\textbf{Herramienta:} Adobe Premiere Pro 2021

\textbf{Efecto utilizado:} ``Cambiar color'' (Corrección de color)

\textbf{Configuración:}
\begin{itemize}
    \item Color objetivo: Verde del semáforo (seleccionado con gotero)
    \item Tolerancia coincidente: 15-30\%
    \item Parámetro variable: Tono (hue rotation)
\end{itemize}

\textbf{Videos generados:}

\begin{table}[H]
\centering
\begin{tabular}{|c|c|c|}
\hline
\textbf{Video} & \textbf{Tono} & \textbf{Descripción} \\
\hline
1 & 0 & Verde original (sin modificación) \\
2 & -20 & Verde levemente modificado \\
3 & -65 & Verde-amarillento \\
4 & -80 & Amarillo-verdoso \\
5 & -100 & Amarillo \\
6 & -110 & Ámbar intenso \\
\hline
\end{tabular}
\caption{Videos generados con variaciones de tono}
\end{table}

\subsection{Resultados}

\textbf{Videos 1-5 (Tono 0 a -100):}

\begin{itemize}
    \item El sistema clasificó el semáforo como GREEN en el 100\% de los frames de la fase verde
    \item Ninguna confusión con YELLOW o RED
    \item El sistema es robusto a estas modificaciones de color
\end{itemize}

\textbf{Video 6 (Tono -110) - Punto de falla:}

\begin{table}[H]
\centering
\begin{tabular}{|l|c|c|l|}
\hline
\textbf{Fase} & \textbf{Frames} & \textbf{Clasificación} & \textbf{Observación} \\
\hline
Fase 1 & 0-165 & GREEN (100\%) & Clasificación correcta \\
\hline
Fase 2 & 166-188 & YELLOW (100\%) & \textbf{Punto de falla} \\
\hline
Fase 3 & 189-443 & RED ($\sim$95\%) & Confusión severa \\
\hline
Fase 4 & 444-484 & BLACK/RED & Inestabilidad \\
\hline
\end{tabular}
\caption{Distribución de clasificaciones en Video 6}
\end{table}

\subsection{Análisis}

\textbf{Hallazgo principal:} El sistema es \textbf{sorprendentemente robusto} a variaciones de color.

\textbf{Umbral de falla:}
\begin{itemize}
    \item El sistema funciona correctamente hasta Tono -100
    \item Falla a partir de Tono -110
    \item El umbral exacto está entre -100 y -110
\end{itemize}

\textbf{Modo de falla:}
\begin{itemize}
    \item La transición es gradual: GREEN → YELLOW → RED
    \item No hay un salto abrupto, sino una degradación progresiva
    \item Frame 166 marca la primera detección errónea (YELLOW)
\end{itemize}

\section{Validación de la Hipótesis}

El test de transición verde→amarillo valida la hipótesis de que el conocimiento interno permite diseñar tests efectivos:

\begin{enumerate}
    \item \textbf{El conocimiento de las clases de salida} (GREEN, YELLOW) sugirió testear la frontera entre ellas

    \item \textbf{El conocimiento de la arquitectura CNN} sugirió que existiría un límite de decisión

    \item \textbf{El conocimiento del proceso de normalización} confirmó que el sistema opera en espacio de color similar al original

    \item \textbf{El resultado} encontró el umbral exacto de falla del clasificador
\end{enumerate}

Un test de caja negra podría haber probado variaciones aleatorias de color sin encontrar sistemáticamente este umbral.

\section{Limitaciones del Test}

\subsection{Adversario Sintético}

El test utiliza un \textbf{adversario sintético}: la modificación de color se realizó digitalmente con un editor de video.

\textbf{Implicación:} No sabemos si esta modificación de color corresponde a algún escenario físico real (iluminación específica, tipo de semáforo, condiciones atmosféricas).

\textbf{Defensa del enfoque:}
\begin{itemize}
    \item El objetivo era encontrar el límite del clasificador, no replicar escenarios físicos
    \item El test es reproducible y sistemático
    \item Los hallazgos sobre la robustez del sistema son válidos independientemente de si el escenario es físicamente realizable
\end{itemize}

\subsection{Cobertura Limitada}

El test solo explora la frontera verde→amarillo. Otros límites de decisión (amarillo→rojo, verde→black) no fueron testeados sistemáticamente.

\section{Justificación del Recorte para Testing}

El test de transición verde→amarillo demuestra el valor del recorte:

\begin{enumerate}
    \item \textbf{Sin el recorte:} Habría sido necesario configurar Apollo completo (Docker, CyberRT, HD-Map) para poder ejecutar este test

    \item \textbf{Con el recorte:} Pudimos ejecutar el test rápidamente, iterar sobre las variaciones de color, y obtener resultados inmediatos

    \item \textbf{Conocimiento interno:} El proceso de reimplementación nos dio el conocimiento necesario para diseñar el test
\end{enumerate}

Esto valida la tesis central del trabajo: \textbf{si focalizamos en un subsistema y lo comprendemos profundamente, podemos diseñar tests más efectivos}.

\section{Resumen}

Este capítulo presentó:

\begin{enumerate}
    \item El enfoque de testing de caja blanca basado en conocimiento interno

    \item El diseño e implementación del test de transición verde→amarillo

    \item Los resultados: umbral de falla entre Tono -100 y -110

    \item La validación de la hipótesis: el conocimiento interno permite encontrar adversarios

    \item Las limitaciones: el adversario es sintético, no necesariamente físicamente realizable
\end{enumerate}